\title{Controlling Text-to-Image Diffusion by Orthogonal Finetuning}

\begin{abstract}
State-of-the-art text-to-image diffusion models demonstrate remarkable proficiency in synthesizing photorealistic images conditioned on textual inputs. However, devising methods for reliably steering or customizing these models to address a variety of downstream objectives remains a significant challenge. In response, we propose Orthogonal Finetuning~(OFT), a theoretically grounded finetuning strategy tailored to adapt text-to-image diffusion models for downstream applications. Distinct from existing approaches, OFT is theoretically shown to maintain hyperspherical energy, a property that captures the pairwise relationships among neurons projected onto the unit hypersphere. Preservation of this hyperspherical energy proves vital for sustaining the semantic generative capacity of diffusion models in text-to-image settings. To further enhance stability during finetuning, we introduce Constrained Orthogonal Finetuning (COFT), which supplements OFT by enforcing a radius constraint on the hypersphere. We focus on two primary text-to-image finetuning scenarios: subject-driven generation—where the goal is to synthesize images featuring a specific subject, based on a limited set of example images and textual guidance—and controllable generation, which involves guiding image synthesis with auxiliary control signals. Empirical evaluation demonstrates that the OFT framework yields superior performance in both image quality and convergence rate compared to alternative methods.
\end{abstract}

\section{Introduction}

Cutting-edge text-to-image diffusion models~\cite{saharia2022photorealistic,ramesh2022hierarchical,rombach2022high} have achieved outstanding capabilities in generating visually compelling images conditioned on text descriptions. Nevertheless, reliance on text prompts alone frequently results in ambiguous or imprecise control, falling short of enabling detailed or specific manipulation of image content. To mitigate these limitations, this paper addresses two specialized categories of text-to-image generation tasks:

\begin{itemize}[leftmargin=*,nosep]

    \item \textbf{Subject-driven generation}~\cite{ruiz2023dreambooth}: When provided with a handful of reference images depicting a target subject, the objective is to synthesize new images featuring that same subject situated in diverse contexts, guided by an accompanying text prompt.
    \item \textbf{Controllable generation}~\cite{zhang2023adding,mou2023t2i}: This task requires the model to generate images that adhere to a supplemental control input (\emph{e.g.}, canny edge detections, segmentation annotations) in conjunction with textual instructions.
\end{itemize}

Both challenges fundamentally center around the problem of how to adapt text-to-image diffusion models to new tasks through finetuning, while maintaining the strong generative capabilities established during pretraining. The following criteria outline what is sought in an effective finetuning approach: (1) \emph{training efficiency}, defined by the need for a minimal set of trainable parameters and reduced training epochs, and (2) \emph{generalizability preservation}, where the goal is to maintain the original model’s capacity to produce high-quality and varied outputs. Conventionally, finetuning is carried out by either slightly modifying the network’s weights using a low learning rate (\emph{e.g.}, \cite{ruiz2023dreambooth}), or by integrating an auxiliary, re-parameterized module with a modest number of additional weights (\emph{e.g.}, \cite{hulora2022,zhang2023adding}). However, both of these straightforward methods are insufficient for ensuring that the generative performance achieved during pretraining remains intact. Moreover, there is a notable gap in principled guidelines for crafting optimal finetuning techniques or for determining essential hyperparameters like the appropriate number of epochs. One of the main obstacles is the absence of an established metric for measuring how well the pretrained generative capabilities are retained. Prevailing finetuning techniques often rely on the heuristic that minimizing the Euclidean distance between the adapted and the original model corresponds to better preservation of pretrained abilities; consequently, small learning rates are the norm. Although this strategy may sometimes yield satisfactory results, we contend that Euclidean proximity does not adequately reflect the true extent of semantic retention. Thus, devising a more structured metric to represent the divergence between finetuned and original models could substantially enhance both the preservation of generative performance from pretraining and the robustness of the finetuning process.

Building on empirical findings that hyperspherical similarity effectively captures semantic structure~\cite{liu2017deep,liu2018decoupled,chen2020angular}, we leverage hyperspherical energy~\cite{liu2018learning} to describe the pairwise relationships among neurons. Hyperspherical energy, computed as the aggregate hyperspherical similarity (\emph{e.g.}, cosine similarity) across all neuron pairs within a layer, quantifies the extent to which neuron representations are uniformly spread over the unit hypersphere~\cite{liu2021learning}. Our working assumption is that optimal finetuning should maintain a hyperspherical energy profile closely aligned with that of the pretrained model. Although simply introducing a regularizer to maintain constant hyperspherical energy during finetuning is a straightforward approach, it does not reliably ensure a minimal difference between the finetuned and original energies. To address this, we exploit an invariance property of hyperspherical energy—namely, that pairwise hyperspherical similarity remains unchanged when all neurons undergo an identical orthogonal transformation. This invariance underpins our proposed method, Orthogonal Finetuning~(OFT), which adapts large-scale text-to-image diffusion models for specific downstream tasks while preserving their hyperspherical energy. The core mechanism involves learning an orthogonal transformation shared across layers that conserves the pairwise angular relationships among neurons. In this sense, OFT can be interpreted as redefining the coordinate system for neurons within a layer. Furthermore, acknowledging that reduced Euclidean distance between finetuned and pretrained parameters is correlated with better retention of pretraining capabilities, we introduce Constrained Orthogonal Finetuning~(COFT), an extension of OFT. COFT further restricts the finetuned parameters to remain within a hypersphere of fixed radius centered at the pretrained neuron locations.

The rationale for employing orthogonal transformation in the context of finetuning neurons is partly inspired by the 2D Fourier transform, a technique that decomposes an image into magnitude and phase spectra. The phase spectrum, encapsulating the angular relationship between the input and the basis functions, is known to retain the bulk of the image’s semantic content. For instance, even if an image’s phase spectrum is paired with an arbitrary magnitude spectrum, the reconstructed image typically maintains its semantic integrity. This observation implies that manipulating neuron orientations is central to inducing semantic changes in generated images—a primary objective of both subject-driven and controllable image generation. Nonetheless, unconstrained modification of neuron directions can severely compromise the generative performance achieved during pretraining. To address this, we advocate for preserving inter-neuronal angles, grounded in the supposition that these angular relationships encode essential neural network knowledge. Consequently, we introduce a layer-wise, shared orthogonal transformation across neurons that maintains constant hyperspherical energy, ensuring angular relationships are unchanged.

Further, our methodology is influenced by work on orthogonal over-parameterized training~\cite{liu2021orthogonal}, where classification neural networks are trained from the ground up through orthogonal transformations applied to randomly initialized parameters. Random initialization is known to yield an expectedly low hyperspherical energy, and the main objective in \cite{liu2021orthogonal} is to preserve this low energy throughout training—an approach linked to improved generalization in classification tasks~\cite{liu2018learning,lin2020regularizing}. Their results show that orthogonal transformations suffice to develop classification networks with strong generalizability. By contrast, our approach centers on finetuning text-to-image diffusion models, targeting enhanced controllability and improved generative capabilities in downstream tasks. It is important to distinguish OFT from the method of \cite{liu2021orthogonal} on two fronts. Firstly, whereas \cite{liu2021orthogonal} is tailored to minimize hyperspherical energy, OFT is explicitly designed to conserve the hyperspherical energy characteristic of the pretrained model, thus preserving its intrinsic structural properties during finetuning—a critical requirement for diffusion models, as reducing hyperspherical energy may compromise semantic structure. Secondly, OFT explicitly limits deviation from the pretrained model, giving rise to a constrained variant, whereas no such restriction is imposed in \cite{liu2021orthogonal}. The challenge in finetuning is to balance adaptability and stability, and we contend that the OFT framework strikes this balance effectively. Our key contributions are summarized as follows:

\begin{itemize}[leftmargin=*,nosep]

    \item We introduce Orthogonal Finetuning, a new approach for finetuning text-to-image diffusion models to enhance their controllability. Additionally, to bolster stability, we develop a constrained version that restricts the angular deviation relative to the pretrained weights.
    \item Distinct from other finetuning techniques, OFT maintains the hyperspherical energy during model adaptation, and we provide empirical evidence that this property is crucial for conserving the generative semantics of the original model.
    \item We implement OFT on two specific applications: subject-driven generation and controllable generation. Through extensive experimental analysis, we show that OFT significantly outperforms earlier methods regarding quality of generated images, finetuning stability, and speed of convergence. OFT is also more data-efficient, achieving robust convergence with fewer training samples and epochs.
    \item For the controllable generation scenario, we propose a novel control consistency metric to assess the degree of control. This metric works by inferring the control signal from generated outputs and then measuring its alignment with the intended control input, thus quantitatively supporting OFT’s enhanced controllability.
\end{itemize}

\section{Related Work}

\textbf{Text-to-image diffusion models}. There has been remarkable advancement~\cite{saharia2022photorealistic,ramesh2022hierarchical,rombach2022high,gu2022vector,nichol2022glide} in generating images from text, largely propelled by innovations in diffusion-based generative approaches~\cite{ho2020denoising,song2021score,song2021denoising,dhariwal2021diffusion} as well as advances in vision-language representation learning~\cite{su2020vlbert,lu2019vilbert,radford2021learning,wang2021simvlm,li2022blip,li2023blip,alayrac2022flamingo,schuhmann2022laion}. In the pixel space, models like GLIDE~\cite{nichol2022glide} and Imagen~\cite{saharia2022photorealistic} have been developed: GLIDE employs a jointly trained text encoder and diffusion prior with paired data, whereas Imagen relies on a pretrained, fixed text encoder. Diffusion models working in latent space include Stable Diffusion~\cite{rombach2022high} and DALL-E2~\cite{ramesh2022hierarchical}: Stable Diffusion leverages VQ-GAN~\cite{esser2021taming} for constructing a visual latent codebook, while DALL-E2 incorporates CLIP~\cite{radford2021learning} to generate a shared latent space for image and text representation. Beyond diffusion models, other generative paradigms such as generative adversarial networks~\cite{reed2016generative,zhang2017stackgan,xu2018attngan,li2019controllable} and autoregressive architectures~\cite{ramesh2021zero,ding2021cogview,wu2022nuwa,yuscaling} have also been explored for text-to-image synthesis. OFT functions as a model-agnostic finetuning technique, making it applicable across various text-to-image diffusion frameworks.

\textbf{Subject-driven generation}. Approaches such as \cite{avrahami2022blended,nichol2022glide} introduce user-provided masks as additional conditions to protect subject features during image editing. Inversion-based strategies~\cite{choi2021ilvr,dhariwal2021diffusion,ramesh2022hierarchical,gal2022image} allow changes to the context while leaving the main subject unaltered. Some techniques like \cite{hertz2022prompt} can facilitate both local and global manipulations even without mask inputs. Nevertheless, these methods often fall short in preserving the fine details linked to subject identity. Pivotal Tuning~\cite{roich2022pivotal} adapts a generator in the vicinity of an inverted latent code, introducing an extra regularization term to retain subject identity. Analogously, \cite{nitzan2022mystyle} derives a tailored generative face prior from a person-specific image collection. \cite{casanova2021instance} supports varied instantiations of an object but can sometimes neglect instance-specific features. DreamBooth~\cite{ruiz2023dreambooth}, utilizing a custom token and a handful of subject images, finetunes text-to-image diffusion models via a reconstruction objective coupled with a class-aware prior preservation loss. Our method, OFT, builds on the DreamBooth scheme; however, rather than simple low-rate finetuning, we utilize orthogonal transformations for model finetuning

\textbf{Controllable generation}. Image-to-image translation represents a subset of controllable generation, where most prior approaches leverage conditional generative adversarial networks~\cite{isola2017image,zhu2017toward,wang2018high,park2019semantic,choi2018stargan}. Recently, diffusion models have also been applied to this task~\cite{saharia2022palette,voynov2022sketch,wang2022pretraining}. Notably, ControlNet~\cite{zhang2023adding} demonstrates the effectiveness of adapting pretrained diffusion models through finetuning with supplementary control inputs, achieving remarkable controllability in generation. Similarly, T2I-Adapter~\cite{mou2023t2i} enhances the controllability of outputs by finetuning pretrained diffusion models. Adhering to the experimental settings introduced in \cite{zhang2023adding,mou2023t2i}, our approach utilizes OFT to finetune pretrained diffusion models, resulting in superior controllability with reduced training data requirements and fewer finetuning parameters. Importantly, OFT incurs no extra computational cost during inference at test time.

\textbf{Model finetuning}. The practice of adapting large pretrained models for downstream applications has seen widespread adoption~\cite{devlin2018bert,brown2020language,he2022masked}. Within this context, a variety of adaptation strategies (\emph{e.g.}, \cite{hulora2022,houlsby2019parameter,pfeiffer2020adapterfusion}) have been comprehensively explored in natural language processing. Among these, LoRA~\cite{hulora2022}, which presumes a low-rank formulation for the additive parameter modifications during finetuning, shares the closest affinity with OFT. In distinction, OFT employs a layer-wise orthogonal transformation that multiplicatively modifies neuron weights and, by design, maintains the pairwise angular relationships between neurons in a layer, contributing to improved training stability.

\section{Orthogonal Finetuning}

\subsection{Why Does Orthogonal Transformation Make Sense?}

We begin by exploring the motivation for employing orthogonal transformation in the finetuning of text-to-image diffusion models. To clarify this, we consider two underlying questions: (1) what is the rationale for adjusting the angular direction of neuron weights, and (2) what motivates the use of orthogonal transformations for modulating these directions.

To address the first question, we take cues from prior empirical findings \cite{liu2018decoupled,chen2020angular}, which suggest that differences in angular features effectively capture the semantic gap. Work such as SphereNet~\cite{liu2017deep} demonstrates that constraining all neurons to lie on a unit hypersphere during neural network training preserves representational power and even enhances generalization. This observation implies that neuron direction encodes essential data features. To further illustrate the significance of neuron angles, we perform a controlled experiment depicted in Figure~\ref{fig:toy_ae}: we train a conventional convolutional autoencoder with flower images. During training, feature maps are computed using the usual inner product between kernel output $z$ (resulting from applying convolution kernel $\bm{w}$ to input $\bm{x}$ in a local window). At test time, we generate feature maps through three methods: (a) the same inner product from training, (b) extraction of only the feature vector magnitudes, and (c) utilization of angular (directional) components only. As shown in Figure~\ref{fig:toy_ae}, angular features nearly perfectly reconstruct input images, while magnitude alone carries negligible semantic content. Notably, we never use the cosine (angle-based) operation during training; only the inner product is used. These outcomes underscore that neuron directions primarily encode the semantic content of input images, suggesting that modifying neuron angles (directions) will be most effective for semantic manipulation through finetuning.

Addressing the second question, an intuitive approach to refining neuron directions is to apply coordinated rotation or reflection to all neurons within the same layer, which equates to orthogonal transformations. Although more general angular adjustments—such as rotating each neuron by distinct angles—are possible, we observe that orthogonal transformations strike an optimal balance between expressiveness and regularity. Supporting this perspective, \cite{liu2021orthogonal} demonstrates that orthogonal transformations are potent tools for neural network training. We reinforce this argument with an experiment conducted in the context of controllable image generation using the ControlNet~\cite{zhang2023adding} setup. The outcomes, illustrated in Figure~\ref{fig:ortho}, reveal that our default OFT achieves reliable and precise control after training converges (by epoch 20). In contrast, an OFT variant omitting the orthogonality constraint completely fails to generate plausible images or achieve controllability. This evidence underscores the critical role of enforcing orthogonality in OFT.

Traditional finetuning methods generally employ gradient descent with a reduced learning rate to update all or selected layers of a model. Using a small learning rate implicitly restricts the updates, thereby limiting the departure from the pretrained parameter values. In effect, standard finetuning seeks to optimize model performance while implicitly keeping $\thickmuskip=2mu \medmuskip=2mu \| \bm{M} - \bm{M}^0\|$ small, where $\bm{M}$ and $\bm{M}^0$ represent the weights of the finetuned and pretrained models, respectively. While this implicit restriction is intuitively justified, it may offer excessive flexibility when applied to very large models. To tighten this, LoRA introduces a direct low-rank constraint on the parameter updates, that is, $\thickmuskip=2mu \medmuskip=2mu \text{rank}(\bm{M} - \bm{M}^0)=r'$, with $r'$ being a predetermined small integer. In contrast, OFT enforces a constraint on the similarity structure among neurons, expressed as $\thickmuskip=2mu \medmuskip=2mu \| \text{HE}(\bm{M})-\text{HE}(\bm{M}^0) \|=0$, where $\text{HE}(\cdot)$ quantifies the hyperspherical energy associated with a given weight matrix.

To clarify, consider a fully connected layer characterized by $\thickmuskip=2mu \medmuskip=2mu \bm{W}=\{\bm{w}_1,\cdots,\bm{w}_n\}\in\mathbb{R}^{d\times n}$, where each neuron $\bm{w}_i\in\mathbb{R}^{d}$ ($\bm{W}^0$ denotes the pretrained weight matrix). For an input $\thickmuskip=2mu \medmuskip=2mu \bm{x}\in\mathbb{R}^d$, the layer produces the output $\thickmuskip=2mu \medmuskip=2mu \bm{z}=\bm{W}^\top\bm{x}$. The OFT perspective is to minimize the discrepancy in hyperspherical energy between the pretrained and finetuned states:

\begin{equation}\label{eq:oft_principle}
\footnotesize
\min_{\bm{W}} \left\| \text{HE}(\bm{W})-\text{HE}(\bm{W}^0)\right\|~~\Leftrightarrow~~\min_{\bm{W}} \bigg{\|} \sum_{i\neq j} \|\hat{\bm{w}}_i-\hat{\bm{w}}_j\|^{-1}-\sum_{i\neq j} \|\hat{\bm{w}}^0_i-\hat{\bm{w}}^0_j\|^{-1}\bigg{\|}
\end{equation}

where the normalized neuron vectors are defined as $\thickmuskip=2mu \medmuskip=2mu \hat{\bm{w}}_i:=\bm{w}_i/\|\bm{w}_i\|$, and hyperspherical energy for the fully connected layer $\bm{W}$ is given by $\thickmuskip=2mu \medmuskip=2mu \text{HE}(\bm{W}):= \sum_{i\neq j} \|\hat{\bm{w}}_i-\hat{\bm{w}}_j\|^{-1}$. The global minimum of Eq.~\eqref{eq:oft_principle} is achieved when the value is zero. This scenario is attained if $\bm{W}$ and $\bm{W}^0$ differ only by an orthogonal transformation, i.e., $\thickmuskip=2mu \medmuskip=2mu \bm{W}=\bm{R}\bm{W}^0$, with $\thickmuskip=2mu \medmuskip=2mu \bm{R}\in\mathbb{R}^{d\times d}$ as an orthogonal matrix (where determinant values $1$ and $-1$ indicate rotation and reflection, respectively). This is the central concept underlying OFT: finetuning proceeds by optimizing layer-wise orthogonal matrices that systematically transform neurons within each layer. Explicitly, for a given pretrained fully connected layer $\thickmuskip=2mu \medmuskip=2mu\bm{W}^0\in\mathbb{R}^{d

\begin{equation}\label{eq:oft_general}
    \footnotesize
    \bm{z} = \bm{W}^\top\bm{x} = (\bm{R} \cdot \bm{W}^0)^\top\bm{x}, ~~~\text{s.t.}~\bm{R}^\top\bm{R} = \bm{R}\bm{R}^\top = \bm{I}
\end{equation}

Here, $\bm{W}$ represents the weight matrix after OFT-based finetuning, and $\bm{I}$ stands for the identity matrix. Figure~\ref{fig:block} visually depicts OFT. Analogous to the practice of zero initialization in LoRA, OFT must ensure that finetuning proceeds exactly from the initial pretrained parameters $\bm{W}^0$. To accomplish this, we set the initial orthogonal matrix $\bm{R}$ as the identity matrix, which allows the finetuned parameters to begin identically with the pretrained ones. Ensuring the orthogonality constraint on $\bm{R}$ can be managed using differentiable orthogonalization techniques, as outlined in \cite{liu2021orthogonal,lezcano2019cheap}. In the following, we detail approaches to maintain orthogonality efficiently in computation.

\subsection{Efficient Orthogonal Parameterization}\label{sect:eff_ortho}

Traditional orthogonalization processes, such as the Gram-Schmidt procedure, though differentiable, typically entail substantial computational overhead~\cite{liu2021orthogonal}. To promote computational efficiency, we utilize the Cayley parameterization for generating orthogonal matrices. In this formulation, the orthogonal matrix takes the form $\thickmuskip=2mu \medmuskip=2mu \bm{R} = (\bm{I} + \bm{Q})(I - \bm{Q})^{-1}$, with $\bm{Q}$ defined as a skew-symmetric matrix that satisfies $\thickmuskip=2mu \medmuskip=2mu \bm{Q} = -\bm{Q}^\top$. While Cayley parameterization is efficient, it restricts the resulting orthogonal matrices to those with determinant $1$, i.e., within the special orthogonal group. Nevertheless, our empirical observations indicate that this restriction does not negatively influence performance. It is important to note that employing the Cayley transform alone may result in high parameter counts for $\bm{R}$ when the dimension $d$ is large. To mitigate this parameter inefficiency, we propose constructing $\bm{R}$ as a block-diagonal matrix comprising $r$ blocks, as expressed below:

\begin{equation}
    \footnotesize
    \bm{R} = \text{diag}(\bm{R}_1, \bm{R}_2, \cdots, \bm{R}_r) = \begin{bmatrix}
    \bm{R}_1 \in \text{O}(\frac{d}{r}) & & \\
    & \ddots & \\
    & & \bm{R}_r \in \text{O}(\frac{d}{r})
    \end{bmatrix} \in \text{O}(d)
\end{equation}

Here, $\text{O}(d)$ refers to the group of orthogonal matrices of dimension $d$, with $\thickmuskip=2mu \medmuskip=2mu \bm{R}\in\mathbb{R}^{d\times d}$ and each block $\thickmuskip=2mu \medmuskip=2mu \bm{R}_i\in\mathbb{R}^{d/r\times d/r}$, $\forall i$, also orthogonal. In the special case where $\thickmuskip=2mu \medmuskip=2mu r=1$, the block-diagonal construction reduces to an unconstrained orthogonal matrix. A general orthogonal matrix of size $\thickmuskip=2mu \medmuskip=2mu d\times d$ involves $\thickmuskip=2mu \medmuskip=2mu d(d-1)/2$ free parameters, corresponding to a computational complexity of $\mathcal{O}(d^2)$. If the matrix is composed of $r$ orthogonal blocks, the total number of parameters decreases to $\thickmuskip=2mu \medmuskip=2mu d(d/r-1)/2$, resulting in $\thickmuskip=2mu \medmuskip=2mu \mathcal{O}(d^2/r)$ complexity. By optionally enforcing all blocks to be identical, i.e., $\thickmuskip=2mu \medmuskip=2mu\bm{R}_i=\bm{R}_j$ for all $i\neq j$, the parameter count can be further minimized to $\mathcal{O}(d^2/r^2)$. Despite these parameter-saving techniques, the overall matrix $\bm{R}$ retains orthogonality, and thus the hyperspherical structure is fully preserved.

We now compare the parameter efficiency of OFT with that of LoRA. Considering LoRA configured with a low-rank parameter $r'$, its number of trainable parameters is $\thickmuskip=2mu \medmuskip=2mu r'(d+n)$. Assuming $r$ and $r'$ both scale with $d$, such that $\thickmuskip=2mu \medmuskip=2mu r=r'=\alpha d$ for some constant $\thickmuskip=2mu \medmuskip=2mu 0<\alpha\leq 1$, the parameter complexity for LoRA is $\mathcal{O}(d^2+dn)$, whereas for OFT it is $\mathcal{O}(d)$. To provide a concrete comparison, suppose the weight matrix has shape $\thickmuskip=2mu \medmuskip=2mu 128\times 128$. LoRA, with $\thickmuskip=2mu \medmuskip=2mu r'=8$, requires $2,048$ trainable parameters, while OFT needs only $960$ trainable parameters when $\thickmuskip=2mu \medmuskip=2mu r=8$ and no block sharing is employed.

\subsection{Constrained Orthogonal Finetuning}

The adaptability of OFT can be further curtailed by restricting the finetuned model to remain in close proximity to the pretrained parameters. Specifically, Constrained Orthogonal Finetuning (COFT) adopts the following forward computation:

\begin{equation}\label{eq:coft_general}
    \footnotesize
    \bm{z}=\bm{W}^\top\bm{x}=(\bm{R}\cdot\bm{W}^0)^\top\bm{x},~~~\text{s.t.}~\bm{R}^\top\bm{R}=\bm{R}\bm{R}^\top=\bm{I},~\norm{\bm{R}-\bm{I}}\leq \epsilon
\end{equation}

Here, the orthogonality requirement on $\bm{R}$ is accompanied by an additional constraint enforcing that $\bm{R}$ remains within an $\epsilon$-ball centered at the identity matrix. While the orthogonality can be addressed through the Cayley parameterization (see Section~\ref{sect:eff_ortho}), integrating the $\epsilon$-deviation constraint within the Cayley framework is not straightforward. To elucidate this, the Cayley transform can be approximated by a Neumann series expansion as $\thickmuskip=2mu \medmuskip=2mu \bm{R}=(\bm{I}+\bm{Q})(I-\bm{Q})^{-1}\

Because the series is convergent under the operator norm, the restriction $\thickmuskip=2mu \medmuskip=2mu\|\bm{R}-\bm{I}\|\leq\epsilon$ can be equivalently imposed within the Cayley transform, leading to a corresponding constraint $\thickmuskip=2mu \medmuskip=2mu \|\bm{Q}-\bm{0}\|\leq\epsilon'$ for $\bm{Q}$, where $\bm{0}$ denotes the zero matrix and $\epsilon'$ is a distinct tolerance parameter from $\epsilon$. Projected gradient descent allows straightforward enforcement of this new condition on $\bm{Q}$. To ensure the initial orthogonal matrix $\bm{R}$ is the identity, $\bm{Q}$ is initialized as a zero matrix. The COFT method thus unites two clear constraints: limiting hyperspherical energy change and bounding deviation from the pretrained parameters. While most prior finetuning approaches only indirectly impose the latter, COFT explicitly introduces it. Although OFT achieves strong performance, we find that COFT provides even greater finetuning stability thanks to this explicit control over deviation. As illustrated in Figure~\ref{fig:eps_coft}, the value of $\epsilon$ modulates how adaptable COFT is during finetuning. If $\epsilon$ is increased, COFT's behavior approaches that of OFT; if decreased, the output aligns more closely with the original pretrained text-to-image diffusion model.

\subsection{Re-scaled Orthogonal Finetuning}

We present a straightforward modification of OFT, in which a trainable scaling parameter is introduced for each neuron. The intuition is that scaling neurons does not influence hyperspherical energy, since the effect is nullified during normalization. Concretely, the forward computation is: $\thickmuskip=2mu \medmuskip=2mu\bm{z}=(\bm{R}\bm{W}^0\bm{D})^\top\bm{x}$\footnote{\textbf{Errata}: The NeurIPS camera-ready version mistakenly presents the re-scaled OFT forward pass as $\thickmuskip=2mu \medmuskip=2mu\bm{z}=(\bm{D}\bm{R}\bm{W}^0)^\top\bm{x}$. The codebase uses the correct formulation, and thus the outcomes in Appendix~\ref{app:rescaled} are unaffected.}, where $\thickmuskip=2mu \medmuskip=2mu\bm{D}=\text{diag}(s_1,\cdots,s_n)\in\mathbb{R}^{n\times n}$ is a diagonal matrix of positive, learnable entries $s_1, \ldots, s_n$. This modification differs from the baseline OFT forward pass in Eq.~\eqref{eq:oft_general}, which makes only $\bm{R}$ trainable; here, both the scaling matrix $\bm{D}$ and the orthogonal transformation $\bm{R}$ are optimized. The introduction of scaling in OFT provides greater model flexibility with only a marginal increase in parameter count. In our experiments, we employ the original OFT to highlight the impact of orthogonal transformation in isolation, but we observe that re-scaled OFT often yields superior results (see Appendix~\ref{app:rescaled}).

\section{Intriguing Insights and Discussions}

\textbf{OFT’s architectural versatility}. OFT can, in principle, be incorporated into any neural network architecture. For instance, within Transformers, LoRA is commonly adopted to update attention matrices~\cite{hulora2022}. For fair comparison, our experiments restrict OFT-based finetuning to attention weights as well. Beyond dense layers, OFT is readily adaptable to convolutional settings, thanks to the interpretable block-diagonal form of $\bm{R}$, which affords advantages over LoRA. By matching the block count to the number of input channels, each channel’s neurons are independently transformed, similar to depth-wise convolution kernels~\cite{chollet2017xception}. If all blocks are tied, then the same orthogonal transformation applies across channels. Preliminary experiments applying OFT to convolution layers can be found in Appendix~\ref{app:convolution}.

\textbf{Connection to LoRA}. LoRA achieves parameter-efficient finetuning by introducing a low-rank adaptation matrix, thus preserving the majority of information stored in the original pretrained weights by restricting the magnitude of change. In comparison, OFT imposes an orthogonality constraint on the transformation applied to the pretrained weights, ensuring that the update does not compromise the pre-learned knowledge. The forward computation in OFT can be reformulated as $\thickmuskip=2mu \medmuskip=2mu \bm{z}=(\bm{R}\bm{W}^0)^\top\bm{x}=(\bm{W}^0+(\bm{R}-\bm{I})\bm{W}^0)^\top\bm{x}$, with the term $\thickmuskip=2mu \medmuskip=2mu (\bm{R}-\bm{I})\bm{W}^0$ bearing a conceptual resemblance to LoRA’s low-rank modification of the weights. Since the initial weight matrix $\bm{W}^0$ is usually full-rank, OFT similarly yields low-rank updates whenever $\thickmuskip=2mu \medmuskip=2mu \bm{R}-\bm{I}$ is of low rank. Mirroring LoRA's use of a tunable rank $r'$, OFT employs a diagonal block size $r$ to further restrict the quantity of learnable parameters. Notably, LoRA and OFT embody two fundamentally different strategies for parameter efficiency: while LoRA capitalizes on the low-rank property, OFT leverages structured sparsity, namely block-diagonal orthogonal matrices.

\textbf{Why OFT converges faster}. There are two primary intuitions underlying OFT’s accelerated convergence. First, as illustrated in Figure~\ref{fig:toy_ae}, effective semantic modification is achieved predominantly by altering the orientation of neurons—precisely the operation OFT facilitates. Second, the OFT framework can be interpreted as adjusting neurons constrained to move on a smooth hyperspherical manifold, leading to a more tractable and favorable optimization landscape. Such benefits are also supported by empirical results in \cite{liu2021orthogonal}.

\textbf{Why not minimize hyperspherical energy}. A major distinction from \cite{liu2021orthogonal} lies in our decision not to pursue minimization of the hyperspherical energy. In the context of classification, non-redundant neuron representations are beneficial. Achieving minimal hyperspherical energy distributes all neurons evenly on the sphere, but this objective may be counterproductive for finetuning, as it can inadvertently erase valuable knowledge captured during pretraining.

\textbf{Balancing flexibility and regularization in finetuning}. Our analysis reveals a fundamental compromise between model adaptability and stability. While conventional finetuning offers maximum adaptability, it tends to lack robustness, frequently resulting in unstable training dynamics or model degradation. In contrast, OFT—a notably straightforward approach—achieves a well-calibrated equilibrium by maintaining the angles between neuron vectors, thus simultaneously enhancing regularity. Additionally, block-diagonal parameterizations can be interpreted as imposing more stringent regularization constraints on the orthogonal matrices.

\textbf{No inference-time performance penalty}. In distinction to frameworks such as ControlNet, OFT does not introduce any extra computational burden during inference for the adapted model. Specifically, at inference, the learned orthogonal transformation $\bm{R}$ can be pre-multiplied with the pretrained weight matrix $\bm{W}^0$, resulting in an equivalent parameter matrix $\thickmuskip=2mu \medmuskip=2mu \bm{W}=\bm{R}\bm{W}^0$. As a result, the inference throughput matches that of the original pretrained configuration.

\section{Experiments and Results}

\textbf{Experimental setup}. For empirical evaluation, we employ Stable Diffusion v1.5~\cite{rombach2022high} as our foundational text-to-image architecture. In each comparison, output samples from all evaluated techniques are chosen at random. We adopt subject-driven generation practices from DreamBooth~\cite{ruiz2023dreambooth}, and employ procedures from ControlNet~\cite{zhang2023adding} and T2I-Adapter~\cite{mou2023t2i} for controllable generation. To maintain a level playing field with LoRA, OFT or COFT modifications are restricted to the identical network layers where LoRA is utilized. Supplemental experimental configurations and findings are outlined in Appendix~\ref{app:settings}.

\subsection{Subject-driven Generation}

\textbf{Configuration details}. DreamBooth~\cite{ruiz2023dreambooth} and LoRA~\cite{hulora2022} serve as baselines in our assessment. Consistent with DreamBooth, the loss formulation is applied uniformly across all approaches. For both DreamBooth and LoRA, we replicate hyperparameter choices from the original publications to ensure best-practice implementations. Extended results and supplementary analyses are presented in Appendix~\ref{app:settings},\ref{app:coft_oft},\ref{app:more_qualitative},\ref{app:failure}.

\textbf{Finetuning stability and convergence}. Our initial analysis centers on assessing the finetuning stability and convergence rates for DreamBooth, LoRA, OFT, and COFT, as presented in Figure~\ref{fig:teaser} and Figure~\ref{fig:db_convergence}. It is evident that both COFT and OFT facilitate stable finetuning for diffusion models. By the 400th iteration, DreamBooth and the OFT-based methods exhibit strong control, whereas LoRA struggles to maintain the identity of the subject. Progressing to 2000 iterations, DreamBooth exhibits mode collapse, generating degraded outputs, and LoRA fails to depict the yellow shirt, often generating yellow fur instead. Conversely, OFT and COFT consistently provide reliable and controlled image synthesis. These findings support that OFT rapidly converges while retaining stability in subject-centric generation tasks. Importantly, the reduced sensitivity to the finetuning iteration count allows users to avoid the burdensome process of hyperparameter tuning across different subjects. Both OFT and COFT enable the use of larger iteration numbers without the need for fine adjustment. When COFT is used with an appropriately chosen $\epsilon$, both the learning rate and total iterations become largely automatic to select.

\textbf{Quantitative comparison}. Adopting the methodology from \cite{ruiz2023dreambooth}, we perform a systematic quantitative evaluation assessing subject fidelity (using DINO~\cite{caron2021emerging} and CLIP-I~\cite{radford2021learning}), alignment to text prompts (CLIP-T~\cite{radford2021learning}), and image diversity (LPIPS~\cite{zhang2018unreasonable}). In this context, CLIP-I measures the mean pairwise cosine similarity between CLIP features of generated samples and reference images. The DINO metric parallels this approach, substituting ViT S/16 DINO embeddings for CLIP representations. CLIP-T evaluates the average cosine similarity between generated image embeddings and corresponding text prompts. LPIPS is computed as the mean cosine similarity between generated images of the same subject under identical textual conditions. Table~\ref{table:dreambooth_final} illustrates that both COFT and OFT substantially exceed DreamBooth and LoRA in the DINO and CLIP-I metrics, with marginally superior or equivalent performance in terms of prompt fidelity and sample diversity. To account for variance, each method is finetuned using 30 distinct random seeds on the same pretrained network. Collectively, these results demonstrate that OFT consistently delivers improved convergence, stability, and final outcomes relative to competing methods.

\textbf{Qualitative comparison}. To better illustrate the advantages of OFT, we present several randomly selected cases of subject-driven generation in Figure~\ref{fig:dreambooth_final}. In the interest of fairness, every example is produced using the same finetuned model under each respective method, with no additional optimization of text prompts for individual outputs. For each approach, we choose the model exhibiting the highest validation CLIP scores. The visual results in Figure~\ref{fig:dreambooth_final} reveal that OFT and COFT consistently achieve strong semantic fidelity in preserving subjects, whereas LoRA frequently struggles to retain subject characteristics (for instance, LoRA entirely loses the subject identity in the bowl scenario). Additionally, OFT and COFT provide more precise alignment between generated outputs and text prompts, while DreamBooth, despite being reliable at identity preservation, often fails to accurately follow the specified prompts (see, for example, the first row in the bowl instance). This qualitative assessment underscores that the OFT framework attains superior balance between controllability and subject preservation. Furthermore, OFT demonstrates robustness to the choice of iteration count, performing reliably even when the number of iterations is increased—an attribute not shared by either DreamBooth or LoRA. Supplementary visual examples are available in Appendix~\ref{app:more_qualitative}. A human evaluation study in Appendix~\ref{app:human} offers additional support for OFT's effectiveness.

\subsection{Controllable Generation}

\textbf{Settings}. We adopt ControlNet~\cite{zhang2023adding}, T2I-Adapter~\cite{mou2023t2i}, and LoRA~\cite{hulora2022} as baseline models. The main paper addresses three challenging controllable generation tasks: Canny edge to image~(C2I) on the COCO dataset~\cite{lin2014microsoft}, segmentation map to image~(S2I) on the ADE20K dataset~\cite{zhou2017scene}, and landmark to face~(L2F) on the CelebA-HQ dataset~\cite{karras2018progressive,xia2021tedigan}. Each method is employed to finetune Stable Diffusion~(SD) v1.5 for 20 epochs across all three datasets. Additional findings are provided in Appendix~\ref{app:more_qualitative},~\ref{app:more_control_task},~\ref{app:failure}.

\textbf{Convergence}. We assess the convergence behavior of ControlNet, T2I-Adapter, LoRA, and COFT on the L2F benchmark, incorporating both quantitative and qualitative analyses. For our quantitative assessment, the principal metric is the mean $\ell_2$ distance calculated between the controlled facial landmarks and the predicted landmarks. Figure~\ref{fig:face_convergence} illustrates the progression of landmark errors for each model across various epochs of finetuning. The results demonstrate that COFT and OFT converge at a markedly faster rate compared to competing methods. Specifically, LoRA requires 20 epochs to achieve a performance level that OFT reaches by epoch 8. Importantly, OFT and COFT maintain a parameter count comparable to LoRA—considerably fewer than ControlNet—while converging with greater efficiency than other baselines. Furthermore, Figure~\ref{fig:teaser} further attests to OFT’s rapid convergence: the example on the right highlights OFT’s superior data efficiency, showing strong convergence using merely 5\% of the ADE20K training samples, outperforming both ControlNet and LoRA. For the qualitative comparison, we restrict the discussion to OFT, COFT, and ControlNet, since ControlNet's landmark error is closest to those of our approaches. Visual evidence in Figure~\ref{fig:face_convergence_qual} reveals stable and consistent convergence for OFT and COFT, with face pose generations progressively aligning to control landmarks. In contrast, ControlNet displays an abrupt gain in controllability around epoch 8, mirroring the sudden convergence phenomenon discussed in \cite{zhang2023adding}. The steady and predictable convergence trajectory of OFT notably enhances its usability in applied settings, as training with OFT yields smoother dynamics, thereby simplifying hyperparameter optimization.

\textbf{Quantitative comparison}. We propose a metric for assessing control consistency to quantitatively evaluate controllable generation performance. This metric involves extracting the control signal from the generated output and directly contrasting it with the input control signal. Specifically, for the C2I task, we report IoU and F1 score; for the S2I task, we assess mean IoU along with both mean and overall accuracy; for the L2F task, we use the mean $\ell_2$ distance measured between the control and predicted landmarks. Detailed explanations of these metrics can be found in Appendix~\ref{app:settings}. We ensure that each method under comparison is run with optimally selected hyperparameters. As presented in Table~\ref{table:control_gen}, the OFT and COFT methods exhibit notably superior and more precise controllability in comparison to other baselines. In our experiments, adapter-based methods (such as T2I-Adapter and ControlNet) display slower convergence rates and achieve poorer final performance. While LoRA outperforms ControlNet on the S2I task, it underperforms in both the C2I and L2F settings. Our observations indicate that LoRA’s performance potential remains modest despite extensive hyperparameter optimization. In contrast, the OFT framework’s effectiveness appears to continue improving with increases in the number of trainable parameters, suggesting the method has yet to reach its capacity limit. Notably, the quantitative evaluation of controllable generation provided in our study constitutes an important contribution, enabling precise measurement of a model’s control ability—contingent on the accuracy of the employed segmentation or detection models.

\textbf{Qualitative comparison}. We further conduct a qualitative assessment of OFT and COFT in comparison to leading approaches such as ControlNet, T2I-Adapter, and LoRA. As illustrated by the randomly generated samples in Figure~\ref{fig:control_final}, OFT and COFT consistently deliver images of high fidelity and realism, alongside precise controllability. Specifically, in the S2I task, LoRA is unable to produce images that conform to the input segmentation map, whereas ControlNet, OFT, and COFT demonstrate strong controllability over generated content. OFT and COFT, when contrasted with ControlNet, are capable of synthesizing images with richer detail and more coherent geometric composition, all while utilizing a significantly smaller number of parameters. Regarding the C2I task, OFT and COFT both successfully produce semantically relevant images from basic Canny edge inputs, outperforming T2I-Adapter and LoRA, which fall short in this respect. For the L2F task, our approach achieves superior pose adherence in face generation, even in challenging scenarios with complex facial orientations. Across all three evaluated control tasks, OFT and COFT yield images of qualitatively higher standard compared to current state-of-the-art baselines, thus substantiating the controllable generation capability of the OFT framework. To broaden the scope of our qualitative analysis, additional results for each control task are included in Appendix~\ref{app:control_exp}. Furthermore, we provide evidence that OFT effectively generalizes to other control modalities—such as dense pose to human body, sketch to image, and depth to image—in Appendix~\ref{app:more_control_task}.

\section{Concluding Remarks and Open Problems}

Inspired by the critical role that angular relationships between neurons play in determining visual semantics, we introduce a straightforward yet powerful finetuning technique—orthogonal finetuning—for text-to-image diffusion model control. Our approach addresses two major tasks: subject-driven generation and controllable generation. Relative to prevailing approaches, OFT offers enhanced controllability and stability during finetuning, while requiring fewer parameters. Significantly, OFT maintains inference efficiency without incurring additional computational costs, enabling practical deployment.

The development of OFT also highlights several intriguing avenues for future research. First, the orthogonality in OFT is enforced via Cayley parametrization, which necessitates the computation of a matrix inverse—a step that poses scalability concerns. While we have mitigated this constraint through the adoption of block diagonal parametrization, devising more efficient, differentiable solutions for the matrix inverse is an unresolved challenge. Second, OFT exhibits distinct compositional capabilities, as orthogonal matrices from multiple OFT finetuning tasks can be sequentially combined while retaining orthogonality. Investigating whether these composed matrices can concurrently encode the knowledge from all constituent downstream tasks is an open question. Lastly, the parameter efficiency of OFT is currently achieved by exploiting a block diagonal matrix structure, which may introduce certain biases and limit adaptability. Identifying alternative strategies that further improve parameter efficiency while minimizing such biases constitutes a significant and open area for further exploration.

\end{document}